\section{Data}
\label{sec:data}

\subsection{Corpus and chunking}

The corpus consists of 82 Russian business documents: construction estimates
and specifications, food-supply invoices and menus, IT contracts, pharmacy
price lists and security policies. Text is obtained with an OCR pipeline
(\texttt{Tesseract} and a lightweight \texttt{PP-OCR} variant as a fallback
\citep{smith2007tesseract,du2020ppocr}), normalized, and split into windows of
at most 8{,}000 tokens at paragraph boundaries. Within a window, text is
further split into chunks of at most 1{,}200 tokens, with tables handled
separately: a table is emitted as a single chunk of at most 1{,}500 tokens and
eight rows, so that headers and their values stay together. Chunks are
contiguous segments of a window, so adjacent chunks share vocabulary; the
median word-level Jaccard similarity of neighboring chunks is 0.31. This
matters at evaluation time, since a retrieved neighbor is frequently a valid
answer (Section~\ref{sec:queries}).

Every chunk is described by metadata: document identifier, category, chunk
index within the document, table flag, token counts, and a content hash of its
window. Near-duplicate documents and duplicate
(window,~chunk) pairs are removed before labeling; the window hash is later
used as the grouping key for the train/validation split, so that no window
appears on both sides.

\subsection{Context labelling}

Each chunk is labeled by an API teacher
(\texttt{deepseek-v4.1-flash}, reasoning disabled). The prompt contains the
full document window and the chunk, and asks for a short context that situates
the chunk within the document and resolves pronouns, abbreviations and table
headers. Prompt caching is effective here, since all chunks of a window share the
same prefix: the overall cache hit rate is 79.3\%, which brings the cost of
labeling the whole corpus down to USD~0.82.

The prompt explicitly forbids adding facts that are absent from the document,
and the audit checks this. Table~\ref{tab:dataset} summarizes the resulting
dataset. The median chunk is 461 tokens long, while the median generated
context is only 296 characters, so the enrichment is cheap to store.

\begin{table}[t]
\centering
\small
\caption{Dataset statistics. Chunks are segments of overlapping document windows, extracted from 82 Russian business documents. Contexts were labeled by an API teacher (\texttt{deepseek-v4.1-flash}, reasoning off) with 79.3\% prompt-cache hit rate.}
\label{tab:dataset}
\begin{tabular}{lr}
\toprule
Property & Value \\
\midrule
Documents & 82 \\
Chunks (window$\rightarrow$context pairs) & 3134 \\
Median chunk length (tokens) & 461 \\
Median generated context (characters) & 296 \\
Table chunks & 1608 \\
\midrule
\quad construction & 211 \\
\quad food & 496 \\
\quad it & 1879 \\
\quad pharma & 239 \\
\quad security & 309 \\
\midrule
Teacher prompt tokens (total) & 24.31M \\
Teacher labeling cost & \$0.82 \\
\midrule
Train / validation chunks & 2701 / 433 \\
Split unit & \texttt{window\_hash} \\
\bottomrule
\end{tabular}
\end{table}


\subsection{Query set}
\label{sec:queries}

Retrieval is evaluated with 234 queries that carry gold chunks. Candidate
queries are generated by two language models from document windows, paraphrased
and judged by one of the two, filtered automatically for identifier leakage and
trivial overlap, and finally reviewed by hand; 234 of the 276 candidates (85\%)
were accepted, at a cost of USD~0.32. The set covers five categories and four
query types: lexical (62), multi-hop (57), semantic (56) and table-oriented (59).

Because chunks are contiguous segments of overlapping windows, a single
question often has several acceptable answers. The evaluation therefore uses
gold groups: for each query, all chunks with the same document identifier whose
chunk index lies within two positions of the gold chunk, or whose text
overlaps the gold chunk with Jaccard similarity above 0.9, are treated as
correct. Metrics that ignore this grouping underestimate dense retrieval, which
tends to return a neighboring window of the correct text.

A manual spot check of twenty query--chunk pairs found seventeen clearly
correct, two debatable and one mismatch. The query set is synthetic, and we
therefore treat absolute metric values as optimistic by roughly five to fifteen
percent; comparisons between arms are made on the same set and are unaffected
by this bias.

\begin{figure}[t]
\centering
\begin{tikzpicture}[
  font=\footnotesize,
  node distance=5.2mm,
  box/.style={draw, rounded corners=1pt, align=center, inner sep=3pt,
              minimum height=6.2mm, text width=30mm, fill=white},
  dev/.style={box, fill=blue!6},
  api/.style={box, fill=orange!8},
  arr/.style={-{Latex[length=1.6mm]}, semithick},
  phase/.style={draw, dashed, rounded corners=2pt, inner sep=5pt}
]

% ---- left column: ingest
\node[box] (ocr) {Documents\\ \scriptsize (OCR)};
\node[box, below=of ocr] (win) {Windows and chunks\\ \scriptsize 8k / 1.2k tokens};
\node[api, below=of win] (gen) {Context generator\\ \scriptsize teacher (API) or student (0.6B)};
\node[dev, below=of gen] (cc) {Contextual chunk\\ \scriptsize context $\Vert$ chunk};
\node[dev, below=of cc] (emb) {Embedder\\ \scriptsize 0.6B, 1024-d};
\node[dev, below=of emb] (idx) {Index\\ \scriptsize Qdrant};

\draw[arr] (ocr) -- (win);
\draw[arr] (win) -- (gen);
\draw[arr] (gen) -- (cc);
\draw[arr] (cc) -- (emb);
\draw[arr] (emb) -- (idx);

% ---- right column: query
\node[box, right=22mm of ocr] (q) {Query\\ \scriptsize with instruction};
\node[dev, below=of q] (qemb) {Embedder\\ \scriptsize 0.6B};
\node[dev, below=of qemb] (ann) {ANN search\\ \scriptsize top-$k$};
\node[box, below=of ann] (bm) {BM25\\ \scriptsize optional};
\node[dev, below=of bm] (fuse) {RRF fusion\\ \scriptsize optional};
\node[box, below=of fuse] (rer) {Reranker\\ \scriptsize optional, 0.6B};
\node[box, below=of rer] (out) {Ranked chunks};

\draw[arr] (q) -- (qemb);
\draw[arr] (qemb) -- (ann);
\draw[arr] (ann) -- (fuse);
\draw[arr] (bm) -- (fuse);
\draw[arr] (fuse) -- (rer);
\draw[arr] (rer) -- (out);

\begin{scope}[on background layer]
  \node[phase, fit=(ocr) (gen), label={[font=\scriptsize\itshape]above:ingest}] {};
  \node[phase, fit=(cc) (idx), label={[font=\scriptsize\itshape]left:on-device}] {};
  \node[phase, fit=(q) (out), label={[font=\scriptsize\itshape]above:query}] {};
\end{scope}

\end{tikzpicture}
\caption{System overview. Ingest (left) is executed once per document; the
dashed frames separate the API teacher from components that run on the target
device. Query processing (right) is executed per request; BM25 fusion and the
reranker are the two optional stages evaluated in Section~\ref{sec:results}.
Without the two optional stages a query takes about 40~ms at the median.}
\label{fig:pipeline}
\end{figure}
